\documentclass[aspectratio=169,10pt,spanish]{beamer}

\input{../../../_preambulo_beamer.tex}
\input{../../../_comandos_beamer.tex}

\selectlanguage{spanish}

\title{MA1001B - Modelación Estadística para la Toma de Decisiones}
\subtitle{Lección 09: Probabilidad condicional e independencia}
\author{}
\institute{Tecnol\'ogico de Monterrey}
\date{}

\begin{document}

\begin{frame}[plain]
  \vspace{1.2cm}
  {\Large\bfseries MA1001B - Modelación Estadística para la Toma de Decisiones\par}
  \vspace{0.7cm}
  {\large Lección 09: Probabilidad condicional e independencia\par}
  \vspace{0.7cm}
  {\color{TecRojo}\rule{\textwidth}{1.2pt}}
  \vspace{0.8cm}

  Facultad MA1001B\\[0.3cm]
  Periodo\\[0.3cm]
  Tecnol\'ogico de Monterrey
\end{frame}

\begin{frame}{La pregunta de probabilidad condicional}
  \begin{alertblock}{Pregunta guía}
    ¿Cómo cambia el espacio muestral relevante al saber que ocurrió un evento,
    y cuándo ese conocimiento deja otra probabilidad sin cambio?
  \end{alertblock}

  \vspace{0.4em}
  Al final de esta lección, usted puede:
  \begin{itemize}
    \item calcular e interpretar probabilidades condicionales;
    \item explicar por qué invertir una condición cambia la pregunta;
    \item probar independencia con igualdades equivalentes de probabilidad;
    \item actualizar una probabilidad de secuencia al muestrear sin reemplazo.
  \end{itemize}

  \vfill
  \scriptsize Todos los registros de pedidos y procesos operativos son sintéticos.
\end{frame}

\begin{frame}{Condicionar reduce el espacio muestral}
  \small
  Sea $M$ revisión manual y $L$ entrega retrasada.

  \begin{center}
    \begin{tabular}{lrrr}
      \toprule
      \textbf{Proceso de enrutamiento} & \textbf{Retrasado} & \textbf{A tiempo} & \textbf{Total} \\
      \midrule
      Revisión manual & 48 & 72 & 120 \\
      Automatizado & 32 & 248 & 280 \\
      \midrule
      Total & 80 & 320 & 400 \\
      \bottomrule
    \end{tabular}
  \end{center}

  \[
    P(L\mid M)=\frac{P(L\cap M)}{P(M)}=\frac{48}{120}=0.40.
  \]

  La condición $M$ convierte los 120 pedidos de revisión manual en el nuevo
  grupo de referencia, en lugar de los 400 pedidos.
\end{frame}

\begin{frame}{Paso 1: La dirección cambia el denominador}
  \begin{columns}[T]
    \begin{column}{0.36\textwidth}
      \small
      \begin{block}{Salida verificada}
        $P(L\mid M)=0.400000$\\
        $P(L\mid A)=0.114286$\\
        $P(M\mid L)=0.600000$
      \end{block}

      $P(L\mid M)$ pregunta por pedidos retrasados dentro de la revisión
      manual. $P(M\mid L)$ pregunta por revisión manual dentro de los
      retrasados. No tienen que coincidir.
    \end{column}
    \begin{column}{0.64\textwidth}
      \centering
      \includegraphics[width=0.93\textwidth]{../figuras/condicional_01_tasas.png}
    \end{column}
  \end{columns}
\end{frame}

\begin{frame}{Independencia significa que no hay actualización}
  \small
  Los eventos $M$ y $L$ son independientes si se cumple cualquier prueba
  equivalente:
  \[
    P(L\mid M)=P(L),
    \qquad
    P(M\cap L)=P(M)P(L).
  \]

  Para los pedidos sintéticos:
  \[
    0.40\ne0.20,
    \qquad
    0.12\ne(0.30)(0.20)=0.06.
  \]

  \begin{alertblock}{Límite de interpretación}
    La tabla muestra una asociación entre enrutamiento y estado de entrega.
    No muestra que la revisión manual cause entrega retrasada.
  \end{alertblock}
\end{frame}

\begin{frame}{Paso 2: Dos pruebas equivalentes llegan a una decisión}
  \begin{columns}[T]
    \begin{column}{0.34\textwidth}
      \small
      \begin{block}{Comparación verificada}
        Condicional: $0.40\ne0.20$\\
        Producto: $0.12\ne0.06$\\[0.2em]
        Independiente: \textbf{Falso}
      \end{block}

      Basta una igualdad fallida para rechazar independencia. La segunda es
      una comprobación de consistencia.
    \end{column}
    \begin{column}{0.66\textwidth}
      \centering
      \includegraphics[width=0.98\textwidth]{../figuras/condicional_02_independencia.png}
    \end{column}
  \end{columns}
\end{frame}

\begin{frame}{Paso 3: Sin reemplazo se crea dependencia}
  \begin{columns}[T]
    \begin{column}{0.39\textwidth}
      \small
      En un lote de 10 pedidos, 4 están marcados. Para dos selecciones sin
      reemplazo,
      \[
        P(F_1\cap F_2)=\frac{4}{10}\frac{3}{9}=0.133333.
      \]

      La independencia ingenua da $0.4^2=0.160000$. Una simulación con
      semilla 42 de 500{,}000 selecciones da \textbf{0.132938}.

      \textbf{Límite:} las secuencias más largas requieren ramas explícitas.
    \end{column}
    \begin{column}{0.61\textwidth}
      \centering
      \includegraphics[width=\textwidth]{../figuras/condicional_03_secuencia.png}
    \end{column}
  \end{columns}
\end{frame}

\begin{frame}{Verificación de conceptos}
  Use la tabla sintética de pedidos y el lote de inspección.
  \begin{enumerate}
    \item ¿Qué denominador corresponde a $P(L\mid M)$?
    \item ¿Por qué $P(L\mid M)=0.40$ es distinto de $P(M\mid L)=0.60$?
    \item ¿Son $M$ y $L$ independientes? Sustente la decisión con una prueba
      de igualdad.
    \item Tras seleccionar un pedido marcado sin reemplazo, ¿por qué la
      siguiente probabilidad de marcado se vuelve $3/9$?
  \end{enumerate}

  \begin{block}{Conexión con la Lección 10}
    Los diagramas de árbol hacen visibles, rama por rama, las
    probabilidades condicionales que cambian y los espacios muestrales de
    varias etapas.
  \end{block}

  \vfill
  \scriptsize Fuente: OpenStax, \textit{Introductory Business Statistics 2e},
  Capítulo 3, Secciones 3.1--3.3. CC BY-NC-SA 4.0.
\end{frame}

\appendix



\end{document}
